This section defines the new instructions and their elementary
well-posedness properties. The first subsection makes a coherent read one
joint tick whose sampled value is stored in $C$ and used by later ticks,
with stopping replaced by a fixed schedule. The second is needed for the
equivalence in Section~\ref{sec:primitives}: it shows an observed, adaptive
protocol is already described by retaining its measurement copies and
tracing at the end. The third is the heart of the paper: keep is a
controlled XOR $K$ from the machine register to quantum targets, a marked
term denotes $U^\dagger KU$, and a \emph{menu} is the one data-dependent
construction admitted by these structured terms---a fixed finite direct
sum $\sum_i\proj i_D\otimes\mathcal U_{P_i}$, with the selector register
$D$ retained, not a data-dependent mark. Soundness, restoration and
resources for these definitions are Sections~\ref{sec:soundness}--%
\ref{sec:resources}, respectively; equivalence of the nine basic
instructions, which excludes coherent reads and marks, is
Section~\ref{sec:primitives} and can be read before those three sections.
The separate conditional bounded dynamic interpreter is
Section~\ref{sec:tierB}.

\subsection{Coherent reads and a fixed stopping schedule}

\begin{definition}[Read instruction, return, and continuation]\label{def:read}
Primitive 10 is exactly a read-mode tick, at its decoded program position.
Before such a tick, its input value is not made available to $f$ by an
external parameter. Later ordinary ticks use the stored value in $C$.
Bounded fragments share that same $C$ and obey the global cell discipline
of Definition~\ref{def:joint-tick}. Intervening gates are composed in their
actual order, including gates changing an earlier read input.
Run a fragment for its supplied number of ticks. Its ordinary return,
stop or fault status selects idle updates for subsequent ticks, but even
an idle update uses the completed $F_t$. The schedule never stops by
observing a coherent status bit. A separate final timeout tick, included
in a declared timeout horizon $T\geq1$, has its own total update
$f_\tau(c,z)=h_\tau(c)$, where $h_\tau(c)$ is obtained from $c$ by
replacing the status field by $\sigma_{\mathrm{to}}$ if it equals
$\sigma_{\mathrm{act}}$, and otherwise preserving every bit of $c$,
on every configuration code. Its decoder
is identically idle. Thus $S_{T-1}=F_{T-1}^{\tau}$ uses $h_\tau$ in
the snapshot formula and $D_{T-1}=\id$; only $S_0,\ldots,S_{T-2}$
use the ordinary device pair $(d,f)$. The final tick has its own blank
$J_{T-1},B_{T-1}$, counted among the $T$ ticks. It is a total unitary
by the same snapshot inverse. It is an explicitly scheduled snapshot
operation, not an assertion that the device's time-independent $f$
can infer the end of its horizon. Without a declared timeout horizon a
fixed term has only its explicitly listed atoms.

For a specified continuation address $\ell$ whose binary code fits
the program-position field, define $h_\ell(c)$ by
replacing just the program-position bits with the binary code of $\ell$
and preserving every other bit, on \emph{all} codes. A continuation atom
uses a fresh $s$-bit $J$ and the same snapshot formula with $f=h_\ell$.
On blank $J$ it installs $\ell$ and retains the previous configuration.
It is a total permutation, including on nonblank $J$.
A structured source block with that continuation denotes the marked
operation below followed by this explicit atom. Its continuation record
is part of $M$ until all enclosing inverses finish. Cleanliness of the
marked operation does not imply cleanliness of this subsequent update.
\end{definition}

In particular, ordinary code after a read acts as a family of reversible
operations controlled by the read register, in the sense of
Lemma~\ref{lem:read-product}, while the operations preserve its
computational-basis sectors. A gate changing that register removes this
guarantee: the controlled description applies to the preceding segment,
and the read register can be entangled with the machine register.
Theorem~\ref{thm:sA-four} is an explicit four-qubit calculation of the
failure, using copying atoms rather than full-snapshot read ticks.

\begin{lemma}[Preserved read register]\label{lem:read-product}
Consider a fixed tick sequence whose only actions on a factor $X\subseteq Q$
are uses as CNOT controls for reads, and whose other operations preserve
each computational-basis sector of $X$. Then on $X L$, where $L$ is all
remaining retained wires, its unitary has the form
\[
 U=\sum_x\proj x_X\otimes V_x,\qquad
 V_x=V_{x,T-1}\cdots V_{x,0}.
\]
Each $V_{x,t}$ is obtained by substituting $x$ in the actual tick circuit;
it is not an independently supplied oracle. If no gates on other data
occur and all other actions are Boolean, these are permutation unitaries.
\end{lemma}
\begin{proof}
 \step{r2read-sector}{Every tick restricts to a unitary $V_{x,t}$ in
 each fixed-$x$ sector.}
 \begin{proof}\pf\ The hypothesis gives an invariant sector. Its
 finite dimension and the unitary tick imply the restriction maps it
 onto itself. Substitution of $x$ computes this restriction.\end{proof}
 \qedstep
 \begin{proof}\pf\ Multiply the sector sums using
 $\proj x\proj y=0$ for $x\ne y$. Products of permutations are
 permutations. A gate changing $X$ violates the hypothesis and gives
 no such formula for the entire sequence.\end{proof}
\end{proof}

\subsection{Bounded adaptivity and its retained interpretation}

\begin{definition}[Bounded adaptive protocol and outputs]\label{def:adaptive}
A protocol is the device data of Definition~\ref{def:device}, an initial
valid classical configuration $c_{\rm in}$ whose live flags identify
exactly the input list, and a horizon of $T\geq1$ ticks including timeout.
A measurement-outcome sequence is \emph{reachable} from $c_{\rm in}$ if
the successive configurations obtained by starting at $c_{\rm in}$ and
applying $f$ with those sampled bits (zero at non-measurement ticks) are
the configurations at which those samples are taken; the final tick uses
$h_\tau$ as specified below.
For every $c$, the decoder range excludes read mode: an intended read
or any other mode violation must select fault with update $c_{\rm fault}$.
All live-flag requirements of Definition~\ref{def:device} hold for every
code, not just for a selected input. There are no marks. At the first
$T-1$ ticks the classical configuration selects the quantum operation,
then $f$ updates it. The last tick uses exactly the idle decoder and
$h_\tau$ of Definition~\ref{def:read}, on both interpretations.
Measurement of wire $i$ has outcomes $z=0,1$ and
unnormalized update $\rho\mapsto P_{i,z}\rho P_{i,z}$, where
$P_{i,z}=\proj z_i\otimes\id$ on other wires. Its outcome is supplied
to $f$ in sample bit zero, with other sample bits zero.
Allocation selects a never-used zero factor in a pool reserved
in advance; the live flags forbid operations on unallocated factors.
At most $T$ allocations are needed beyond the input list. Terminal data
are padded as above. Designated quantum output wires $O\subseteq Q$ are
fixed, and for every measurement-outcome sequence reachable from
$c_{\rm in}$ every wire in $O$ is live in the final configuration $c_T$;
the classical output is
the total specified bit string $o(c_T)$, which includes any requested
outcome, memory and console fields. All pool wires outside $O$, including
live non-output wires, are traced in this output instrument. These rules
define it by finite repeated composition and summation over outcomes.

To retain the protocol use exactly $S_{T-1}\cdots S_0$, with the last
factor defined by $h_\tau$, on the input state and the reserved zero
non-input pool, records and samples, with $C$ initialized to $c_{\rm in}$.
Copy $o(C)$ into a fresh output register
$Y$ by its total XOR evaluator, and copy $Y$ into a retained register $E_Y$.
The environment is exactly the complement of $OY$ in the entire retained
list: in particular it contains $Q\setminus O$, $C$, all $J_t,B_t$,
$E_Y$, and evaluation workspace. The second copy makes $Y$
classical after that environment is traced. This specifies the output
extraction, rather than merely observing measurement samples.
\end{definition}

\begin{lemma}[Bounded deferred observation, proved here]\label{lem:deferred}
For a protocol satisfying all the global decoder and live-flag restrictions
and the $T\geq1$ timeout convention in Definition~\ref{def:adaptive},
the retained construction there is an isometry.
Tracing its designated environment gives exactly the protocol's output
instrument, including classical outputs, for inputs with any reference.
\end{lemma}
\begin{proof}
 \step{r2defer-expand}{After $t$ ticks the retained vector can be written
 $\sum_{\zeta} A_{t,\zeta}\ket\psi\otimes\ket{e_{t,\zeta}}$, where
 $\zeta$ is the length-$t$ list of sampled bits, padded by zeros at
 nonmeasurement ticks. Distinct possible lists have orthogonal retained
 sample registers. $A_{t,\zeta}$ is the ordered product of the protocol's
 selected gates and measurement projectors on the reserved pool, composed
 with the initial zero embedding of its non-input wires; its configuration is the
 deterministic configuration prescribed by those samples.}
 \begin{proof}
  \step{r2defer-base}{The assertion holds at $t=0$.}
  \begin{proof}\pf\ There is one empty list, the initial zero embedding,
  and the specified initial configuration and blank records.\end{proof}
  \step{r2defer-next}{The assertion at $t$ implies it at $t+1$.}
\begin{proof}
 \step{r2defer-next-hyp}{\assume{The asserted expansion holds at
 $t<T$.}\prove{It holds after the next scheduled tick.}}
 \begin{proof}
  \step{r2defer-next-timeout}{\case{$t=T-1$.}}
  \begin{proof}\pf\ The scheduled quantum action is identity and
  the new sample is zero. The snapshot applies exactly $h_\tau$
  and stores the old configuration, giving the prescribed next
  configuration without altering the quantum product.\end{proof}
  \step{r2defer-next-gate}{\case{$t<T-1$ and the decoder selects
  gate mode.}}
  \begin{proof}\pf\ Multiply each summand by the selected gate.
  The sample is zero; the snapshot computes $f(c,0)$ and stores $c$.
  This is the corresponding protocol step.\end{proof}
  \step{r2defer-next-measure}{\case{$t<T-1$ and the decoder selects
  measurement mode on wire $i$.}}
  \begin{proof}\pf\ The copying isometry is
  $v\mapsto\sum_{z=0}^1P_{i,z}v\otimes\ket z_{B_t}$,
  with unused sample bits zero. The snapshot stores the old
  configuration and computes precisely the next one from $z$.
  These are the two projector terms of the protocol update.\end{proof}
  \step{r2defer-next-allocate}{\case{$t<T-1$ and the decoder selects
  allocation mode.}}
  \begin{proof}\pf\ Its zero live flag implies the factor is a
  non-input factor initially zero and never acted on: ordinary
  updates preserve flags, allocation only changes zero to one,
  and every quantum action requires a one flag. A fault is terminal
  and cannot be followed by allocation. The factor is therefore
  still zero. The identity quantum action and prescribed flag update,
  with zero sample, implement allocation from the reserved pool.\end{proof}
  \step{r2defer-next-other}{\case{$t<T-1$ and the decoder selects
  fault, ordinary or idle mode.}}
  \begin{proof}\pf\ These modes have identity quantum action,
  zero new sample and their specified updates. Their snapshots
  store the old configuration. Read mode is excluded on every code.\end{proof}
  \step{r2defer-next-orthogonal}{Distinct extended sample lists
  have orthogonal retained sample registers.}
  \begin{proof}\pf\ Distinct old lists differ in an old retained
  sample bit. Equal old lists with different new samples differ
  in $B_t$. No case erases these bits.\end{proof}
  \qedstep
  \begin{proof}\pf\ The cases exhaust the schedule and permitted
  modes. Each appends the correct gate or projector and deterministic
  configuration update; \stepref{r2defer-next-orthogonal} establishes
  the required orthogonality. Sum over the old sample lists.\end{proof}
 \end{proof}
 \qedstep
 \begin{proof}\pf\ By \stepref{r2defer-next-hyp}.\end{proof}
\end{proof}
  \qedstep
  \begin{proof}\pf\ Finite induction using
  \stepref{r2defer-base} and \stepref{r2defer-next}.\end{proof}
 \end{proof}
 \step{r2defer-output}{Tracing the environment gives the required instrument.}
 \begin{proof}\pf\ Orthogonality of the samples removes cross terms
 for different lists in \stepref{r2defer-expand}. The extra copy of
 $o(c_T)$ removes off-diagonal classical output labels. The remaining
 terms are exactly $A_{T,\zeta}\rho A_{T,\zeta}^\dagger$ grouped by
 $o(c_T)$, with all $Q\setminus O$ wires traced, as in the definition.\end{proof}
 \qedstep
 \begin{proof}\pf\ The initial embeddings and subsequent unitaries are
 isometries. The calculation in \stepref{r2defer-output} applies to
 matrix units by linearity, and every operator acts as identity on a
 reference. It therefore applies to all density operators with references.
 No unsupplied measurement theorem is used.\end{proof}
\end{proof}

\subsection{Structured marks, allocation types, and a finite menu}

\begin{definition}[Keep]\label{def:keep}
For the specified current machine list $M$ of $m_M$ bits, choose a total
Boolean circuit $v:\{0,1\}^{m_M}\to\{0,1\}^{|r|}$ and a fixed ordered
list $r$ of distinct quantum wires disjoint from $M$. Set
\[
 K_{v,r}=\sum_{u\in\{0,1\}^{m_M}}X_r^{v(u)}\otimes\proj u_M,
 \qquad X_r^z=\bigotimes_i X_{r_i}^{z_i},                 \tag{2}
\]
with identity on all other wires. A keep list is a finite product of
these maps. Different entries may share target wires; repeated XORs
have their usual meaning. Every target is in the register context of
its enclosing mark. There is no computed-address exception in this
definition. Targets used by the body are allowed, without a cleanliness
promise. A standalone use is a circuit atom with the same typing.
\end{definition}

\begin{lemma}[Keep well-posedness]\label{lem:keep}
Every keep list is a unitary involution and has a clean finite-alphabet
implementation.
\end{lemma}
\begin{proof}
 \step{r2keep-sectors}{In a fixed $u$ sector each entry is a product
 of Pauli $X$ operators; all entries commute and square to identity.}
 \begin{proof}\pf\ Every entry preserves $M$, and $X^2=\id$ on each
 target. Overlaps merely add their exponents modulo two.\end{proof}
 \qedstep
 \begin{proof}\pf\ Sum over the orthogonal $u$ sectors for unitarity
 and the involution identity. Evaluate $v$, XOR its output into $r$,
 and uncompute by Lemma~\ref{lem:synthesis}; $M$ is preserved.\end{proof}
\end{proof}

We use \emph{Tier~A} for the class defined next: marks and their pairings are
fixed in the syntax, hence known to the classical control computer before any
quantum input. \emph{Tier~B} is the data-dependent analogue, in which
whether a rewind is reached and which mark it names can depend on values read
from qubits. Section~\ref{sec:tierB} defines a bounded interpreter and proves
conditional results under a supplied certificate; constructing such a
certificate for general dynamic programs remains open
(Sections~\ref{sec:tierB} and \ref{sec:open-problems}).

\begin{definition}[Tier~A, fixed-space terms and typed outer programs]
\label{def:tiera}
Fix a retained list $L=QM$. An atom is an occurrence-labelled elementary
quantum gate from $\mathcal G^{\pm}$ (or the explicitly declared ideal
alphabet), a joint tick of Definition~\ref{def:joint-tick}, a continuation atom, a specified
completed RTM step, or a finite composition of these. Each is a specified
logical unitary, with an exact clean implementation where finite-alphabet
resources are claimed. Boolean elementary gates on machine bits are
also atoms by Lemma~\ref{lem:synthesis}.
Every atom specifies a support list contained in $L$ and acts as identity
on its complement. An RTM atom's support is its specified
$m_{\rm RTM}$-bit encoding register. For a joint-tick atom the supplied
decoder, on \emph{every} $c$, excludes measurement and allocation modes;
mode violations must instead select fault with update $c_{\rm fault}$.
Read mode is allowed, with all the live-flag conditions of
Definition~\ref{def:device}. The explicitly scheduled timeout snapshot
of Definition~\ref{def:read} is also an allowed joint-tick atom.
Atom labels include syntactic occurrence identifiers: two written copies
of a tick receive distinct recording cells even if their instruction
data agree. Adjoint expansion retains the original occurrence identifiers.
The grammar, interpreted on the \emph{same} $L$, is
\[
 P::=I\mid a\mid P;P\mid \mathsf{mark}\{P\}\mathsf{rewind}\ K.
\]
Denote the last constructor by $\mathsf M(P,K)$ and its term denotation
by $\mathcal U_P$. Set
\[
 \mathcal U_I=\id,
 \quad\mathcal U_a=A_a,
 \quad\mathcal U_{P;R}=\mathcal U_R\mathcal U_P,
 \quad\mathcal U_{\mathsf M(P,K)}
   =\mathcal U_P^\dagger K\mathcal U_P.       \tag{3}
\]
The sans-serif $\mathsf M(P,K)$ is a term constructor, not the register $M$.
The mark is a static delimiter in machine 1. Inverting a completed inner
term takes its complete adjoint, including its keep. A source block with
post-rewind address $\ell$ is, by definition,
$\mathsf M(P,K);\mathsf{cont}_\ell$, with the atom of
Definition~\ref{def:read}. The first operation has exactly the denotation
$\mathcal U_P^\dagger K\mathcal U_P$; it need not restore the machine-held
position. Position restoration requires a separately proved condition
on the allowed input space. The second operation installs $\ell$ when
its own fresh record is blank, retaining the previous configuration,
regardless of the position left by the first operation.

For changing register lists the outer grammar is
$E::=I_L\mid A_B\mid\mathsf{run}_L(P)\mid E;E$,
with types $I_L:L\to L$, $A_B:L\to LB$ for globally fresh $B$,
$\mathsf{run}_L(P):L\to L$, and composition only when the middle
lists coincide. Its denotations are respectively identity, zero embedding,
$\mathcal U_P$, and ordered composition. Outputs are a
designated subset; dropping is performed afterward on its complement.
All logical records and implementation workspace are allocated before
their enclosing run. Disjoint recording cells are assigned to syntactic
occurrences, not to the number of gates in the expanded circuit.
There is a padding rule: a term $P$ on $L$ may be used on $LZ$ for any
disjoint list $Z$, by lifting every atom and keep with tensor identity;
its denotation is $\mathcal U_P\otimes\id_Z$. In particular
$\mathsf{run}_{LZ}(P\otimes I_Z)$ is the logical run after allocating
workspace. A physical implementation may use $Z$, and agrees with this
logical run on zero $Z$ by Definition~\ref{def:scratch}; equality on
other $Z$ states is not required.

\emph{Restrictions of this tier:} no allocation, drop, observed measurement,
external console observation, or actual stop inside a marked body. Bounded
idle statuses and updates of an unobserved console tape are allowed.
Every mark instance and pairing is static. A loop or call containing marks
is covered only if it has a supplied fixed finite syntactic expansion;
data-dependent numbers of mark visits and general recursive calls are
excluded. Bounded loops without marks can be joint-tick fragments.
Source allocation inside a mark is not automatically admitted by this
grammar. Pre-reserving its wires is a correction requiring a separate
zero-input justification, not equality on arbitrary wire states.
\end{definition}

\begin{lemma}[Algebraic typing of structured terms]\label{lem:term-types}
Every finite fixed-space term of Definition~\ref{def:tiera} denotes a
total unitary. Every well-typed outer term denotes an isometry into its
entire retained output space. These conclusions concern the algebraic
denotation, without an assertion about compilation cost or source coverage.
\end{lemma}
\begin{proof}
 \step{r2types-induction}{Every fixed-space term denotes a unitary.}
 \begin{proof}
  \step{r2types-base}{\case{The term is identity or an atom.}}
  \begin{proof}\pf\ Identity is unitary. Elementary gates, completed
  permutations and joint ticks are unitary by their matrices,
  Lemmas~\ref{lem:completion} and \ref{lem:joint}. The continuation
  has the snapshot inverse, and finite compositions of atoms multiply
  unitaries.\end{proof}
  \step{r2types-sequence}{\case{The term is $P;R$, and the assertion
  holds for its proper subterms.}}
  \begin{proof}\pf\ The product of the two unitary denotations has
  inverse the reversed product of their adjoints.\end{proof}
  \step{r2types-mark}{\case{The term is $\mathsf M(P,K)$, and the
  assertion holds for $P$.}}
  \begin{proof}\pf\ Put $U=\mathcal U_P$. By
  Lemma~\ref{lem:keep}, $K$ is unitary; multiplication gives
  \[
   (U^\dagger KU)^\dagger(U^\dagger KU)=\id.
  \]
  The opposite product also equals identity.\end{proof}
  \qedstep
  \begin{proof}\pf\ Structural induction on the finite syntax, by
  \stepref{r2types-base}, \stepref{r2types-sequence}, and
  \stepref{r2types-mark}.\end{proof}
 \end{proof}
 \step{r2types-outer}{Every well-typed outer term is an isometry.}
 \begin{proof}\pf\ Identity and run are isometries by
 \stepref{r2types-induction}. Allocation satisfies $A_B^\dagger A_B=\id$
 because $\ket0_B$ has unit norm. If $V^\dagger V=\id$ and
 $W^\dagger W=\id$ with matching middle lists, then
 $(WV)^\dagger WV=V^\dagger V=\id$. Structural induction gives
 the assertion.\end{proof}
 \qedstep
 \begin{proof}\pf\ By \stepref{r2types-induction} and
 \stepref{r2types-outer}.\end{proof}
\end{proof}

\begin{lemma}[Moving an unused allocation earlier]\label{lem:reserve}
Let $B$ be fresh, and let a finite circuit segment $V$ before its
allocation act on $L$ alone. Then
$(V\otimes\id_B)A_B=A_BV$. Consequently finitely many allocations can
be moved to the start of a forward circuit if none of their names is used
before its original allocation. The equality is of embeddings with
$B$ initialized to zero, not of an inverse allocation with a unitary.
\end{lemma}
\begin{proof}
 \step{r2reserve-tensor}{Both sides send $\psi$ to $V\psi\otimes\ket0_B$.}
 \begin{proof}\pf\ Substitute the definition of $A_B$ and the
 tensor-product action on distinct factors.\end{proof}
 \qedstep
 \begin{proof}\pf\ Move each fresh allocation across each earlier
 operation by \stepref{r2reserve-tensor}. There are finitely many such
 pairs. This proves the equality on the stated input space only.
 It does not define an adjoint deletion inside a marked body.\end{proof}
\end{proof}

\begin{definition}[Restricted menu; dynamic source semantics remains open]
\label{def:tierb}
Let $D$ be a $d$-qubit selector disjoint from a common retained payload
$L$. Choose $0\leq j\leq2^d$ and Tier~A terms $P_0,\ldots,P_{j-1}$
on $L$. For the implementation claim, each term must supply an explicit
exact clean circuit over $\mathcal G^{\pm}$, with its actual matrix
phase, on the same zero-initialized workspace list after padding.
An ideal-alphabet term without such a supplied circuit is excluded from
that claim; the algebraic direct sum still exists for any total payload
unitaries. No general exact controlled-rotation synthesis is presumed.
All $D$ wires are preserved throughout this operation and retained for
every enclosing inverse. Define
\[
 \mathsf{SELECT}=
 \sum_{i<j}\proj i_D\otimes\mathcal U_{P_i}+
 \sum_{j\leq i<2^d}\proj i_D\otimes\id_L.                  \tag{4}
\]
Implement this by enumerating the $j$ circuit lists, controlling each
elementary gate on $D=i$. The schedule is fixed; shorter lists may be
padded with identity. No selected circuit acts on $D$. If $D$ was
computed earlier, that computation is an explicit reversible operation
with retained records; no erasure of $D$ is implied by selection.

The finite-menu extension of the Tier~A syntax admits this construction
as a derived atom $\mathsf{sel}_\nu(P_0,\ldots,P_{j-1})$ on $DL$,
with identity on any other named factors. Its support contains $D,L$;
the global machine list includes $D$ and the payload's machine factors.
The entries here are base Tier~A terms, not menus, so this rule is not
circular. A derived atom may occur in a sequence or a marked body.
Every payload tick and continuation has its own globally fresh occurrence
identifier across all entries and menu occurrences. Reversing this atom
reverses its entire controlled circuit with the same identifiers and
selector wires; it introduces no new cells. Its cost as an atom is the
length of that supplied controlled circuit. Its internal payload ticks
and continuations are counted once for logical record allocation.

More generally a supplied list of $N\geq0$ total \emph{unitaries}
$S_t:\cH_L\to\cH_L$ is a finite unitary schedule. Calling a member
``rewind'' supplies no source semantics. Such arbitrary schedules are
not a definition of dynamic mark/rewind and carry no equivalence claim.
The construction admitted by this definition is (4), with its stated
static menu. Theorems~\ref{thm:tB-sound} and \ref{thm:tB-cost} separately
prove bounded dynamic-interpreter soundness and costs under the supplied
certificate of Definition~\ref{def:tB-cert}. A translation of source-level
mutable mark stacks, repeated visits, and data-dependent inverse stopping
into that certified interpreter remains open
(Open problem~\ref{prob:tB-compiler}).
\end{definition}

\begin{lemma}[Menu well-posedness]\label{lem:menu}
Under Definition~\ref{def:tierb}, if the payload terms denote unitaries,
(4) is a unitary on $DL$. For the supplied exact finite-alphabet lists
required there, its specified controlled implementation has the same
logical action, with zero workspace restored. Adding such derived atoms
to Tier~A preserves the conclusions of Lemma~\ref{lem:term-types}.
\end{lemma}
\begin{proof}
 \step{r2menu-blocks}{The summands have disjoint selector sectors
 covering $\cH_D$, and unitary payloads.}
 \begin{proof}\pf\ Capacity $j\leq2^d$ makes every projector defined;
 the binary basis resolves identity. Multiply (4) by its adjoint in
 either order to obtain identity.\end{proof}
 \step{r2menu-implementation}{In sector $i<j$ precisely list $i$ acts,
 with its specified phase; in other sectors every gate is identity.}
 \begin{proof}\pf\ The equality predicates depend only on the
 unchanged $D$. Lemma~\ref{lem:synthesis} controls each elementary
 gate exactly, and Lemma~\ref{lem:scratch} restores workspace. No
 independent scalar phase may be discarded from a selected list.\end{proof}
 \qedstep
 \begin{proof}\pf\ By \stepref{r2menu-blocks} and
 \stepref{r2menu-implementation}, followed by linearity across sectors.
 For the extension, the new atom case is \stepref{r2menu-blocks};
 the product, marked-term and outer-isometry induction steps in
 Lemma~\ref{lem:term-types} are unchanged. For implementations use
 \stepref{r2menu-implementation} and Lemma~\ref{lem:scratch}.
 \end{proof}
\end{proof}

\subsection{Costs, certificates and cleanliness: definitions used later}
The definitions in this subsection are used in Sections~\ref{sec:primitives}--%
\ref{sec:resources}. They are collected here, next to the marked terms they
describe; no result is asserted in this subsection.

\begin{definition}[Tick selection cost]\label{def:selection-cost}
For the data of Definitions~\ref{def:device} and \ref{def:adaptive}, write
\[
\mathcal R=\beta+\Delta+s+k+(q^2+kq+1)(s+k+1).
\]
The supplied sizes include Boolean-circuit input and output ports. Uniform families and unary numerical bounds have the meaning of Definition~\ref{def:rtm}.
\end{definition}

\begin{remark}[Explicit corrections to the primitive list]\label{rem:corrections}
The corrections to the nine instructions that the equivalence requires
are listed, with their justifications, in Section~\ref{sec:summary}; the
additional restrictions for marked code are those of
Definition~\ref{def:tiera} for Tier~A and
Definition~\ref{def:tB-cert} for the conditional Tier~B interpretation.
\end{remark}

\begin{definition}[Expansion and cost model]\label{def:costs}
Assign every atom and keep an explicit clean circuit over $\mathcal G^{\pm}$
of length $c_a,c_K$. One gate in this alphabet costs one sequential gate
tick (a gate tick, the time unit, is not a joint tick $S_t$; one joint
tick costs many gate ticks), with no parallel execution assumed. Define the unsimplified expansion
by concatenation and by body, keep, reversed adjoint body for a mark.
Let $C(P)$ be its gate count, $L(P)$ the sum of all $c_a,c_K$ over
syntactic occurrences counted once, and $d(P)$ the maximum number of
enclosing marks: zero on atoms and identity, maximum under sequence,
and one plus body depth at a mark. A zero-initialized scratch pool of
width $a_*$ is the maximum local implementation workspace.

Let $n_{\rm tick}(P)$ count distinct syntactic joint-tick occurrences,
$n_{\rm cont}(P)$ continuation occurrences, and $q_{\rm new}(P)$ declared
fresh data wires in the enclosing typed outer program. With one shared
$s$-bit configuration and $m_{\rm extra}$ explicitly declared other
logical wires (including any RTM or Boolean-value registers), the derived
logical machine size for the snapshot construction is
\[
 m_M=s+(s+k)n_{\rm tick}(P)+s n_{\rm cont}(P)+m_{\rm extra}.
\]
For the menu extension, each derived menu atom is one leaf of cost
equal to its compiled controlled-list length in $C(P),L(P),d(P)$.
Its payload occurrences contribute to $n_{\rm tick},n_{\rm cont}$ for
record allocation, but their gate costs are not counted again outside
the atom cost. For selectors of widths $\delta_\nu$, add
$\sum_\nu\delta_\nu$ fresh disjoint selector wires to
the displayed machine size unless already included in $m_{\rm extra}$,
and count the records of their explicit computations once. The same
selector names are retained under adjoints. Selectors computed temporarily during
one tick are clean scratch and are not retained logical selectors.
Compiler time means steps of a multitape classical TM, with explicit
output gate list and binary wire labels; it is distinct from gate ticks.
\end{definition}


\begin{definition}[Approximation certificates]\label{def:certificates}
For a fixed expanded list of $N_g\ge1$ unitary factors $A_i$, a certificate procedure supplies, for every rational $0<\eta<1$, a clean $\mathcal G^{\pm}$ implementation of a logical unitary $\widetilde A_i$ with
\[
\|\widetilde A_i-A_i\|\le\eta,\qquad
\text{length}\le c_*\log^\alpha(2+1/\eta),
\]
where $c_*>0$, $\alpha\ge0$ are fixed independently of $i,N_g,\eta$.
The error concerns actual matrices, including the entire controlled matrix when present. Inverses use adjoints of the same chosen approximating circuits. Occurrences inside inverses and controls are counted in $N_g$.
\end{definition}

\begin{remark}[Scope of the approximation hypothesis]\label{rem:sk}
The certificate hypothesis above is an explicit restriction: no
unconditional approximation theorem for arbitrary additional gates is
proved here. Its removal for arbitrary effectively specified bounded-arity
gates over $\mathcal G$, with all control phases accounted for, remains
an open synthesis task.
Dawson--Nielsen's instruction-set hypothesis is a finite inverse-closed
subset of $SU(d)$ whose generated subgroup is dense in $SU(d)$; their Theorem~1 gives
operator-norm length $O(\log^c(1/\eta))$ for fixed $d$
\cite[Definition~1, PDF p.~2; Theorem~1, PDF p.~4]{DawsonNielsen2006}.
This statement is not applied literally to $\{H,T,\operatorname{CNOT}\}$:
its determinants and inverse closure do not satisfy those hypotheses.
The exact controlled finite gates needed above were instead constructed
in Lemma~\ref{lem:synthesis}. No approximation step relies on an
unread source or on discarding a control-dependent phase.
\end{remark}

\begin{definition}[Cleanliness, reached subspaces, and restoration]
\label{def:clean}
For a fixed unit vector $\ket{m_0}_M$, define
$E\psi=\psi\otimes\ket{m_0}$ and $P_0=EE^\dagger$.
A unitary $W$ is clean on this input space if $WE=EV$ for a unitary
$V$ on $Q$. For a particular input density operator the weaker condition
is support of its final state in $\operatorname{ran}E$.
For nested use let $\mathcal S\subseteq\cH_L$ be a specified reached
subspace with orthogonal projection $P_{\mathcal S}$. It is computed
mathematically as the image of the initial allowed subspace under the
fully specified preceding isometry, or it can be a supplied containing
subspace. It need not factor across data and machine registers.
Given a designated reusable auxiliary factor $A\subseteq L$, put
$P_A=\id_{L\setminus A}\otimes\proj0_A$. Reset on $\mathcal S$ means
$(\id-P_A)WP_{\mathcal S}=0$; for reuse the incoming subspace must
also satisfy $\mathcal S\subseteq\operatorname{ran}P_A$.
For an initially entangled machine factor, preservation of its reduced
state is a different condition, called marginal restoration, and does
not authorize replacing it by fresh zeros or tracing it before an inverse.

Cleanliness here concerns the logical machine register $M$. It is a
different property from a \emph{clean implementation}
(Definition~\ref{def:scratch}), which concerns only implementation
workspace $Z$: a marked operation can have a clean implementation while
leaving $M$ entangled with $Q$ (Theorem~\ref{thm:sA-four}), and
Theorem~\ref{thm:sA-sound} asserts the former kind of statement, not
cleanliness of $M$. In phrases such as ``clean XOR'' or ``clean CNOT'',
clean means that the auxiliary factors of the stated initialized
computation return to their initial vectors; it is not merely a
workspace assertion.
\end{definition}


\begin{definition}[Restricted delayed-erasure comparison]\label{def:delayed}
For $1\leq j\leq d$ let
$f_j:\{0,1\}^{b_{j-1}}\to\{0,1\}^{b_j}$ be deterministic.
A specified total unitary evaluator $A_j$ of cost $a_j$ preserves its
input and, with output and workspace zero, writes
$(f_j(x),g_j(x))$ in computational basis into an output of $b_j$ bits
and private workspace of $u_j$ bits. It touches no other factors.
Copying that output into a disjoint external target costs $k_j$;
compute, copy, inverse gives a composite XOR unitary of cost
$o_j=2a_j+k_j$. The eager composite at level $j$ computes the previous
composite into a zero intermediate, applies the level-$j$ XOR unitary,
and inverts the previous composite. The delayed construction evaluates $A_1,\ldots,A_d$ on
distinct retained intermediates and workspaces, copies the last value,
then inverts $A_d,\ldots,A_1$. Targets are never evaluator inputs and no
other gate changes the retained values during these sequences.
Theorem~\ref{thm:resA-delayed} uses these compositions under a more general
evaluator hypothesis, (RA8) there, of which the basis-valued $g_j(x)$ above
is the special case $\ket{\gamma_{j,x}}=\ket{g_j(x)}$.
\end{definition}


\begin{remark}[Open point: remaining scope]\label{problem:remaining}
Theorems~\ref{thm:tB-sound}--\ref{thm:tB-cost} give conditional bounded
interpreter results, a static specialization, and the exact disposal
criterion. They do not construct a certified interpreter for general
dynamic source mark/rewind, data-dependent repeated mark visits or an
inverse stopped at a coherent mark identifier
(Open problem~\ref{prob:tB-compiler}). The full-snapshot form of the
four-qubit example also remains open.
They also do not settle general phase-correct approximation certificates
for arbitrary additional gates or an optimal recursive resource bound.
Source allocation within a marked body requires a proved reservation
translation and remains excluded from Tier~A as stated.
The proved deferred-observation lemma concerns bounded adaptive protocols
without marks. It does not authorize observed measurement inside a
reversible marked body.
These open points carry no proof step in the settled results.
An optimal time--space trade-off for general intervening quantum gates,
including a transformation to delayed erasure, also remains open.
\end{remark}

\begin{remark}[Collected restrictions]\label{rem:r3-restrictions}
The complete restrictions are collected in Section~\ref{sec:summary}.
They include the supplied-circuit and final-instrument scope of the
nine-instruction equivalence, the static Tier~A rules, and the conditional
Tier~B certificate and implementation hypotheses; none is inferred from
an instruction name alone.
\end{remark}
